% source: c5be9ff6e69851a10537e5802c232ac84b479a3b
\chapter{线性代数基础}
\label{ch:linalg}

\section{本章为何重要}
本章并不是一门通用的线性代数课程。它的目的是以足以在后文派上用场的精确程度，阐述后量子密码赖以建立的那些特定结构：有限域上的向量空间、线性映射及其矩阵，以及决定线性方程组何时可解的秩--零化度关系。格密码的核心对象是方程
\[
A\mathbf{s} + \mathbf{e} = \mathbf{b},
\]
它是一个被小误差 $\mathbf{e}$ 故意扰动的线性方程组。要理解为什么未扰动的方程组容易求解，而扰动后的方程组（据推测）是困难的，我们首先需要精确地说明“求解线性方程组”是什么意思。

\section{学习目标}
读完本章后，你应当能够：

\begin{itemize}
  \item 陈述向量空间的公理，并认识到 $\mathbb{F}_q^n$ 是我们的工作范例；
  \item 定义线性无关、张成、基和维数，并知道为什么维数是良定义的；
  \item 将矩阵描述为线性映射的坐标形式，并定义其秩、核与像；
  \item 陈述秩--零化度定理，并用它来推理 $A\mathbf{x}=\mathbf{b}$ 的可解性；
  \item 精确解释为什么 $A\mathbf{s}=\mathbf{b}$ 是容易的，而 $A\mathbf{s}+\mathbf{e}=\mathbf{b}$ 却是困难的 LWE 问题。
\end{itemize}

\section{标量：底层的域}
每个向量空间都建立在一组可以进行加、减、乘、除运算的\emph{标量}之上。这组标量构成一个\emph{域}（见第~\ref{ch:sagemath}~章，那里介绍了有限域 $\mathbb{F}_q = \mathrm{GF}(q)$）。由于本章其余部分都依赖于这一定义，我们在此重述一遍。

\begin{definition}[域]
\emph{域}是一个带有两种运算 $+$ 和 $\cdot$ 的集合 $\mathbb{F}$，使得 $(\mathbb{F},+)$ 是以 $0$ 为单位元的阿贝尔群，非零元素 $(\mathbb{F}\setminus\{0\},\cdot)$ 构成以 $1$ 为单位元的阿贝尔群，并且乘法对加法满足分配律。特别地，每个非零元素都有乘法逆元。
\end{definition}

我们熟悉的域有 $\mathbb{Q}$、$\mathbb{R}$ 和 $\mathbb{C}$。密码学则是在有限域 $\mathbb{F}_q$ 上进行运算；当 $q$ 为素数时，它就是模 $q$ 的整数，
\[
\mathbb{F}_q = \mathbb{Z}/q\mathbb{Z} = \{0,1,\dots,q-1\},
\]
其算术运算均模 $q$ 约简。使它成为域而不仅仅是环的性质在于可以做除法：由于 $q$ 是素数，每个非零剩余都有逆元。正是这一点使高斯消元得以进行，这也是为什么诸如 $q=3329$（Kyber）这样的密码参数要选为素数。

\section{向量与向量空间}
在中学里，向量是一个坐标元组，例如 $(1,2)$。就我们的目的而言，向量是\emph{向量空间}中的一个元素：一种由域元素构成、可以相加和数乘的结构化容器。

\begin{definition}[向量空间]
域 $\mathbb{F}$ 上的\emph{向量空间}是一个集合 $V$，它带有加法 $V\times V\to V$ 和数乘 $\mathbb{F}\times V\to V$，使得 $(V,+)$ 是阿贝尔群，并且对所有 $a,b\in\mathbb{F}$ 及 $\mathbf{u},\mathbf{v}\in V$，有
\[
a(\mathbf{u}+\mathbf{v}) = a\mathbf{u}+a\mathbf{v},\quad
(a+b)\mathbf{v} = a\mathbf{v}+b\mathbf{v},\quad
(ab)\mathbf{v} = a(b\mathbf{v}),\quad
1\mathbf{v} = \mathbf{v}.
\]
\end{definition}

\begin{example}[空间 $\mathbb{F}_q^n$]
列向量的集合
\[
\mathbb{F}_q^n = \left\{ (v_1,\dots,v_n)^\top : v_i \in \mathbb{F}_q \right\},
\]
在逐分量加法和数乘下，构成 $\mathbb{F}_q$ 上的向量空间。它是有限的：恰好包含 $q^n$ 个元素。Kyber 的秘密钥位于 $\mathbb{F}_{3329}^{256}$ 中，因此它是 $3329^{256}$ 个可能向量之一——这已经远远超过可观测宇宙中的原子数目。
\end{example}

一个典型的秘密写作
\[
\mathbf{s} = \begin{bmatrix}2\\5\\1\\7\end{bmatrix} \in \mathbb{F}_{q}^{4},
\qquad\text{or in practice}\qquad
\mathbf{s} \in \mathbb{F}_{3329}^{256}.
\]

\subsection{Sage 实验：创建向量}

\begin{lstlisting}[language=Python]
F = GF(17)
v = vector(F, [1, 2, 3])
v
\end{lstlisting}

要查看它的类型，可使用 \texttt{type(v)}。Sage 报告它是一个\emph{自由模元素}，而非普通列表。这并非吹毛求疵：域上的向量空间是环上\emph{模}的特例，即该环恰好是一个域的情形；而模 LWE（第~\ref{ch:mlwe}~章）所用的正是这种更一般的对象。

\subsection{向量加法与数乘}

\begin{lstlisting}[language=Python]
a = vector(F, [1, 2, 3])
b = vector(F, [5, 6, 7])
a + b        # (6, 8, 10)
3 * a        # (3, 6, 9)
\end{lstlisting}

所有算术运算都模域的特征进行：在 $\mathbb{F}_7$ 上，和 $6+5$ 等于 $4$ 而非 $11$。在 $\mathbb{R}$ 或 $\mathbb{Q}$ 上，加法具有图~\ref{fig:vector-addition}所示的熟悉的平行四边形图景。

\begin{figure}[H]
\centering
\begin{tikzpicture}[scale=0.9]
  \draw[pqcgray!40, thin] (-0.6,0) -- (4.8,0);
  \draw[pqcgray!40, thin] (0,-0.6) -- (0,3.6);
  \draw[pqcgray!40, dashed] (3,1) -- (4,3);
  \draw[pqcgray!40, dashed] (1,2) -- (4,3);
  \draw[basisvec] (0,0) -- (3,1);
  \draw[altvec] (0,0) -- (1,2);
  \draw[shortvec] (0,0) -- (4,3);
  \node[pqcblue, font=\small] at (3.35,0.7) {$u$};
  \node[pqcteal, font=\small] at (0.55,2.3) {$v$};
  \node[pqcorange, font=\small] at (4.35,2.85) {$u+v$};
\end{tikzpicture}
\caption{按平行四边形法则进行的向量加法，其中 $u=(3,1)$，$v=(1,2)$。两个向量都从原点画出；补全平行四边形（虚线边）可以看出，和 $u+v=(4,3)$ 恰好是它的对角线。按坐标来看，这与上面的 Sage 实验逐分量计算的规则相同：$(1,2,3)+(5,6,7)=(6,8,10)$。}
\label{fig:vector-addition}
\end{figure}

\section{线性组合、线性无关与张成}
整个学科中最重要的单一运算，就是构造\emph{线性组合}。

\begin{definition}[线性组合与张成]
给定向量 $\mathbf{v}_1,\dots,\mathbf{v}_k \in V$ 和标量 $c_1,\dots,c_k\in\mathbb{F}$，向量 $\sum_{i=1}^k c_i\mathbf{v}_i$ 称为这些 $\mathbf{v}_i$ 的一个\emph{线性组合}。所有这样的组合构成的集合称为它们的\emph{张成}，记作 $\mathrm{span}(\mathbf{v}_1,\dots,\mathbf{v}_k)$；它是 $V$ 中包含每个 $\mathbf{v}_i$ 的最小子空间。
\end{definition}

\begin{definition}[线性无关]
若等于零向量的线性组合只有平凡组合，则称向量 $\mathbf{v}_1,\dots,\mathbf{v}_k$ \emph{线性无关}：
\[
\sum_{i=1}^k c_i \mathbf{v}_i = \mathbf{0} \quad\Longrightarrow\quad c_1 = \cdots = c_k = 0.
\]
否则称它们\emph{线性相关}，这意味着其中某一个是其余向量的线性组合。
\end{definition}

例如，$(1,0)$ 与 $(0,1)$ 线性无关，而 $(1,2)$ 与 $(2,4)$ 线性相关，因为后者是前者的两倍。

\subsection{Sage 实验：通过秩检测线性相关}
线性相关可以借助\emph{秩}（下文给出精确定义）在数值上检测：将这些向量作为行堆叠起来，秩就统计了其中真正线性无关的向量个数。

\begin{lstlisting}[language=Python]
A = Matrix(GF(17), [[1, 2], [2, 4]])
A.rank()      # 1: only one independent direction
\end{lstlisting}

\section{基与维数}
\begin{definition}[基]
向量空间 $V$ 的一组\emph{基}是一个既线性无关又张成 $V$ 的向量集合。等价地，每个 $\mathbf{v}\in V$ 都可以\emph{唯一地}表示为基向量的线性组合。
\end{definition}

$\mathbb{F}_q^2$ 的标准基是 $\{(1,0),(0,1)\}$，因为 $(a,b)=a(1,0)+b(0,1)$。但 $\{(2,1),(1,1)\}$ 同样是该空间的一组合法的基：空间是内在的，而基是为它选取的一个坐标系。基中向量的个数则不是可以选择的。

\begin{theorem}[维数是良定义的]
有限维向量空间 $V$ 的任意两组基所含元素个数相同。这个公共的数称为\emph{维数} $\dim V$。
\end{theorem}

因此 $\dim \mathbb{F}_q^n = n$。可以在不改变空间的情况下自由更换基，这正是整个格基约化理论的萌芽（第~\ref{ch:lll}~章）：格同样有许多组基，而找到一组\emph{好的}基，是核心的算法问题。

\subsection{Sage 实验：同一空间的两组基}

\begin{lstlisting}[language=Python]
B1 = Matrix(GF(17), [[1, 0], [0, 1]])
B2 = Matrix(GF(17), [[2, 1], [1, 1]])
# both have rank 2, so each is a basis of GF(17)^2
print(B1.rank(), B2.rank())
\end{lstlisting}

\section{作为线性映射的矩阵}
理解矩阵的最佳方式，不是把它看作一张数字网格，而是把它看作\emph{线性映射}的坐标形式。

\begin{definition}[线性映射]
若映射 $T\colon \mathbb{F}^n \to \mathbb{F}^m$ 保持两种运算，则称它是\emph{线性的}：
\[
T(\mathbf{u}+\mathbf{v}) = T(\mathbf{u})+T(\mathbf{v}),\qquad T(a\mathbf{v}) = a\,T(\mathbf{v}).
\]
一旦取定了基，每个线性映射都由一个 $m\times n$ 矩阵 $A$ 通过 $T(\mathbf{v})=A\mathbf{v}$ 给出，而每个这样的矩阵也都定义了一个线性映射。
\end{definition}

乘积 $A\mathbf{v}$ 本身就是一个线性组合：它是以 $\mathbf{v}$ 的各分量为权重、对 $A$ 的各\emph{列}所作的组合。具体地说，结果的第 $i$ 个分量是 $A$ 的第 $i$ 行与 $\mathbf{v}$ 的内积，这恰好就是算法~\ref{alg:matvec}中的双重循环。

\begin{lstlisting}[language=Python]
F = GF(17)
A = Matrix(F, [[1, 2], [3, 4]])
v = vector(F, [5, 6])
A * v
\end{lstlisting}

\begin{algorithm}[H]
\caption{矩阵--向量乘积}
\label{alg:matvec}
\begin{algorithmic}[1]
\Require 矩阵 $A \in \mathbb{F}^{m\times n}$，其元素为 $A_{i,j}$；向量 $v \in \mathbb{F}^n$
\Ensure $w = Av \in \mathbb{F}^m$
\For{$i = 1,\dots,m$}
    \State $w_i \gets 0$
    \For{$j = 1,\dots,n$}
        \State $w_i \gets w_i + A_{i,j}\, v_j$
    \EndFor
\EndFor
\State \Return $w$
\end{algorithmic}
\end{algorithm}

这同一个循环是全书的基础：它计算 LWE 方程中的 $A\mathbf{s}$；而在环的形式下（第~\ref{ch:rlwe}~章），它又以 RLWE 和 Kyber 所使用的教科书式多项式乘法的面貌重新出现。

\section{秩、核与像}
有三个子空间刻画了线性映射 $A\colon\mathbb{F}^n\to\mathbb{F}^m$ 的作用。

\begin{definition}[像、核、秩]
\emph{像}（或列空间）是 $\operatorname{im}A = \{A\mathbf{v} : \mathbf{v}\in\mathbb{F}^n\}\subseteq\mathbb{F}^m$；其维数称为 $A$ 的\emph{秩}。\emph{核}（或零空间）是 $\ker A = \{\mathbf{v}\in\mathbb{F}^n : A\mathbf{v}=\mathbf{0}\}\subseteq\mathbb{F}^n$；其维数称为\emph{零化度}。
\end{definition}

像恰好是使 $A\mathbf{x}=\mathbf{b}$ 有解的所有右端项 $\mathbf{b}$ 构成的集合。核则衡量解偏离唯一性的程度：若 $A\mathbf{x}_0=\mathbf{b}$，则全部解构成的集合为 $\mathbf{x}_0 + \ker A$。这两个维数是相互关联的。

\begin{theorem}[秩--零化度定理]
对任意线性映射 $A\colon \mathbb{F}^n \to \mathbb{F}^m$，有
\[
\dim(\ker A) + \operatorname{rank}(A) = n.
\]
\end{theorem}

特别地，方阵 $A\in\mathbb{F}^{n\times n}$ 可逆，当且仅当其秩为 $n$，等价地 $\ker A=\{\mathbf{0}\}$，又等价地其各列构成 $\mathbb{F}^n$ 的一组基。当 $A$ 可逆时，对每个 $\mathbf{b}$，$A\mathbf{x}=\mathbf{b}$ 都恰有一个解。

\subsection{Sage 实验：秩决定可逆性}

\begin{lstlisting}[language=Python]
F = GF(17)
for _ in range(10):
    A = random_matrix(F, 4, 4)
    print(A.rank())
\end{lstlisting}

对于 $\mathbb{F}_q$ 上均匀随机的 $n\times n$ 矩阵，其可逆的概率为
\[
\prod_{i=1}^{n}(1-q^{-i}).
\]
因此上述论断依赖于 $q$：对于这个 $\mathbb{F}_{17}$ 上的 $4\times4$ 示例，该概率很高，而在 $\mathbb{F}_2$ 上则只有约 $30.8\%$。

\section{求解线性方程组，以及向 LWE 的跨越}
给定 $A$ 和 $\mathbf{b}$，方程 $A\mathbf{x}=\mathbf{b}$ 可用\emph{高斯消元}求解：通过行变换将 $A$ 化为阶梯形，再经回代求出 $\mathbf{x}$。

\begin{proposition}[线性方程组是容易的]
在域上，高斯消元使用 $O(n^3)$ 次域运算即可求解 $n\times n$ 方程组 $A\mathbf{x}=\mathbf{b}$（或报告无解）。若 $A$ 可逆，则解唯一。
\end{proposition}

正是这种高效性，使得普通的线性方程组\emph{毫无}密码学价值：任何看到 $A$ 和 $\mathbf{b}=A\mathbf{s}$ 的人都能在立方时间内恢复秘密 $\mathbf{s}$。LWE 改变了这一情形。给定一个公开的、通常为长方形的随机矩阵 $A\in\mathbb{Z}_q^{m\times n}$（其中 $m>n$）、一个从规定的小分布中抽取的秘密 $\mathbf{s}$，以及一个独立的小误差向量 $\mathbf{e}$，人们观察到的样本为
\[
\mathbf{b} = A\mathbf{s} + \mathbf{e},
\]
其算术运算模 $q$ 进行。与可逆的方阵方程组不同，对于非零的随机误差，这个超定方程组一般没有精确解。高斯消元可以检测出不相容性，但无法从含噪样本中恢复植入的秘密。恢复 $\mathbf{s}$ 就是\emph{带误差学习}（Learning With Errors）问题（第~\ref{ch:lwe}~章）。

\begin{remark}
值得停下来体会一下：纸面上的变化如此之小，难度上的变化却如此之大。对于可逆方阵，$\mathbf{s}\mapsto A\mathbf{s}$ 很容易求逆。而在有许多含噪样本时，观察到的值还依赖于独立采样的误差；密码学上的困难性就在于区分和利用这种结构。
\end{remark}

\section{实验}

\begin{lstlisting}[language=Python]
F = GF(17)
# 1. random vectors
V = VectorSpace(F, 4)
for _ in range(5): print(V.random_element())

# 2. a matrix acting on a vector
A = random_matrix(F, 6, 4); v = random_vector(F, 4)
print(A * v)

# 3. exact solve is easy; try recovering s from b = A*s
s = random_vector(F, 4); b = A * s
print(A.solve_right(b) == s)   # True when A has full column rank

# 4. perturb an overdetermined system: it normally becomes inconsistent
e = vector(F, [1, 0, 0, 16, 0, 0])   # 16 == -1 mod 17
try:
    A.solve_right(b + e)
except ValueError:
    print("no exact solution: noisy samples are inconsistent")
\end{lstlisting}

该实验只是一个玩具版的示意：真正的 LWE 会精确规定维数以及秘密和误差的分布。

\section{与下一章的联系}
下一章\emph{多项式环}保留了线性代数的视角，但把“标量与向量”的图景改换为一个更丰富的单一对象。一个结构化的（循环）矩阵将由一个多项式来编码，矩阵--向量乘法将变为多项式乘法，而向量空间 $\mathbb{F}_q^n$ 将升级为商环
\[
R_q = \mathbb{Z}_q[x]/(x^n + 1),
\]
这正是 RLWE 和 Kyber 的天然归宿。

\printbibliography[heading=chapterrefs]
